\documentstyle[%


righttag%


,11pt%


,amssymb%


,russian%       remove in Eng.


,izvru%         Set izven% in Eng.


]{amsart}





% {,}


% change BMPs, esp. pic. 7 (...-versions)


\newcommand{\ri}{{\mathrm{I}}}


\newcommand{\rii}{{\mathrm{II}}}


\newcommand{\riii}{{\mathrm{III}}}


\newcommand{\Span}{\operatorname{span}}


\newcommand{\diag}{\operatorname{diag}}


\newcommand{\cond}{\operatorname{cond}}


\newcommand{\tts}[1]{}


\newcommand{\W}{\overset{0}{W}{}^1_2(\Omega)}


% \newcommand{\bc }[1]{\mbox{\boldmath{$\mathcal {#1}$}}}


\newcommand{\bc}[1]{\boldsymbol{\cal {#1}}}


% \newcommand{\bca}{\boldsymbol{\cal A}}


% \newcommand{\bi}[1]{\mbox{\boldmath $ {#1} $}} % Bold italic type


\newcommand{\bi}[1]{\boldsymbol {#1}}


\newcommand{\mc}[1]{\cal {#1}}


\newcommand{\mbb}[1]{\bold {#1}}


\numberwithin{equation}{section}


%\renewcommand{\baselinestretch}{0.97}





\theoremstyle{plain}


\theorembodyfont{\it}


\newtheorem{thms}{\hskip\parindent Теорема}[section]


\newtheorem{lems}{\hskip\parindent Лемма}[section]


\newtheorem{props}{\hskip\parindent Предложение}[section]





\theoremstyle{definition}


\theorembodyfont{\rm}


\newtheorem{cors}{\hskip\parindent Следствие}[section]


\newtheorem{notes}{\hskip\parindent Замечание}[section]





%\hoffset=-15mm% For Eng.


\setcounter{page}{10}


\god{2005}


\nomer{1~(512)} %   Remove in Eng.


%\nomer{Vol.\,49, No.\,1}%    Set in Eng.


%\str{\pageref{begin}--\pageref{end}}%   Set in Eng.


\udk{517.929}





\author{\it И.Е.\,Ануфриев, В.Г.\,Корнеев}


% NB: ^^ remove in Eng.


\title{Численное тестирование метода декомпозиции\\ для $p$-версии МКЭ}


\thanks{Работа выполнена при финансовой поддержке


Российского фонда фундаментальных исследований


(проект \No\,00-01-00772), программы ``Университеты России'' (проект


\No\,УР.03.01.017) и гранта для молодых кандидатов наук Санкт-Петербурга


(\No\,PD03-1.1-22).}





\date{}


%\pagestyle{myheadings}  %         Set in Eng.


\pagestyle{plain} %              Remove in Eng.


\begin{document}


%\thispagestyle{plain}%          Set in Eng.


\maketitle


\label{begin}














\section*{Введение}


\label{intro} 





Известно, что применение $hp$-версии метода конечных элементов для 


решения эллиптических краевых задач дает ряд преимуществ по сравнению 


с $h$-версией. В частности, скорость ее сходимости как правило выше, а 


в ряде важных для приложений случаев она обладает экспоненциальной 


скоростью сходимости \cite{BS1987}, \cite{Schwab1998}. Это позволяет


сократить число неизвестных при решении сложных инженерных задач. В 


то же время матрицы жесткости $hp$-версии существенно более заполнены 


и имеют более сложную структуру. Поэтому в полной мере ее


достоинства могут быть реализованы только при наличии быстрых алгоритмов 


решения соответствующих


систем линейных алгебраических уравнений.





Предположим, $\bold {Av}={\bold f}$ --- симметричная


положительно определенная (с.\,п.\,о.) система алгебраических уравнений с


$n\times n$-матрицей, и $n$ --- параметр. Алгоритм решения этой


задачи назовем {\em быстрым} или (асимптотически)


{\em почти оптимальным}, если он требует $ O((1+\log n)^kn)$


арифметических операций с фиксированным и


небольшим $k$. Если $k=0$, то алгоритм называется {\em оптимальным}. 


Для решения часто применяется предобусловленный метод сопряженных 


градиентов (ПМСГ) с некоторым предобусловливателем $\bc A$.


Под (общей) {\em арифметической или вычислительной стоимостью


предобусловливателя} понимается арифметическая стоимость


решения системы ПМСГ за исключением умножений матрицы ${\bc A}$ на


векторы, неизбежных при итерационном решении. Иными словами,


эффективность предобусловливателя измеряется не только


относительной обусловленностью, но и стоимостью решения системы


с матрицей в виде предобусловливателя. Соответственно общей арифметической 


стоимости предобусловливатель называется быстрым, почти оптимальным и


оптимальным. Для краткости алгоритмы решения систем линейных


алгебраических уравнений называются также {\em солверами}.


В этой работе рассматриваются только иерархические


$hp$-дискретизации эллиптических уравнений второго порядка. Их


конечные элементы предполагаются ассоциироваными с квадратными


базисными элементами, координатные функции которых определяются


посредством произведений интегрированных полиномов Лежандра.


Поэтому по умолчанию всегда подразумевается этот тип


дискретизаций.





Активное использование предобусловливателей и солверов МДО $hp$-версии 


началось в\linebreak 90-х гг., когда для $h$-дискретизаций такие алгоритмы были


уже достаточно хорошо изучены (см., напр., 


\cite{Nepomnyaschikh:91} (1991), \cite{SmithBjorstadGropp:96} (1996)


и ссылки в них). По


иерархическим $hp$-дискретизациям в 1994--97~гг. был выполнен ряд работ, 


в которых рассматривались вопросы, связанные c блочно-диагональными 


предобусловливателями, в которых блок неизвестных в вершинах


элементов независим от остальных, с предобусловливанием дополнений 


Шура, с двухуровневыми МДО и МДО с перекрытием областей и т.п. (см. 


\cite{BabuskaCraigMandelPitkaranta:91}--\cite{guo}).


Наиболее разработанная схема построения МДО-предобусловливателя 


типа Дирихле--Дирихле для двумерных эллиптических задач была представлена


в \cite{KI1995}--\cite{KI1996} и 


была развита затем в \cite{KJ1997}--\cite{KJ1999a}.


Основные достижения этих работ заключались 


в трех моментах. Один --- обоснование возможности предобусловливания с 


потерей $(1+\log p)$ в относительной обусловленности матрицей из двух 


независимых блоков, соответствующих неизвестным в вершинах элементов и 


остальным неизвестным (см. также \cite{BabuskaCraigMandelPitkaranta:91}). 


Два других --- получениe эффективных предобусловливателей


конечно-разностного вида для внутренних задач на подобластях декомпозиции и


предобусловливателей-солверов для интерфейсных задач. Наиболее трудоемкой 


компонентой алгоритмов МДО для $hp$-версии является решение локальных 


задач Дирихле на подобластях декомпозиции. Хотя использование полученного 


конечно-разностного предобусловливателя уменьшало вычислительную


стоимость этой компоненты, построить на его основе почти оптимальные 


солверы для локальных задач Дирихле удалось лишь в более поздних работах 


\cite{K2001}, \cite{K2002}. Они представляют 


собой вторичные итерационные процессы МДО для конечно-разностной задачи,


соответствующей предобусловливателю. В целом эти результаты привели к 


нескольким предобусловливателям-солверам МДО


для иерархических $hp$-дискретизаций эллиптических уравнений


второго порядка, имеющим арифметическую стоимость 


$O(p^2 (1+\ln p)^k)$ c $1\le k \le 3{,}5$. Они представлены


в работах \cite{Korneev:02b}, \cite{KXA2002}.





В данной работе рассматривается одна из версий почти


оптимального МДО алгоритма, положенного в основу


разрабатываемого нами пакета программ. Основным ее содержанием


является описание алгоритмов пакета и результатов их анализа и обсуждение 


результатов достаточно обширных численных экспериментов. Основной целью


тестирования МДО солвера в целом и всех его основных компонент


было на данном этапе установление асимптотик по $p$ их практической


трудоемкости и сравнение их с априорными оценками. Насколько 


известно, в работе впервые представлены численные результаты,


полученные с помощью комплекса программ МДО на базе $hp$-версии.


Проделан ряд экспериментов со всеми


компонентами МДО-предобусловливателя с целью подтверждения


полученных ранее в~\cite{KJ1999}, \cite{KJ1999a},


\cite{KXA2002} априорных оценок


спектральной эквивалентности и объема вычислительной работы. Как


известно, основными компонентами являются следующие.





1. Солверы для локальных задач Дирихле на подобластях 


декомпозиции --- конечных элементах.





2. Процедуры продолжения и ограничения с межэлементной


границы внутрь конечных элементов.





3. Предобусловливатель-солвер для дополнения Шура,


получаемого после конденсации внутренних для конечных элементов


неизвестных.





4. Солвер для решения систем с матрицей, являющейся блоком матрицы 


жесткости, соответствующим


неизвестным в вершинах конечных элементов.





Заметим, что ни конденсация внутренних неизвестных, ни дополнение


Шура в алгоритме не используются. Мы не рассматриваем алгоритмы


для четвертой компоненты, поскольку указанный блок является


матрицей жесткости для ансамбля конечных элементов первого


порядка. Для систем с такими матрицами можно использовать целый


набор быстрых алгоритмов, полученных для $h$-версии.





Тестирование производилось для весьма широкого диапазона


значений $p$, в некоторых экспериментах до $p>500$, выходящих за


пределы, характерные в настоящее время для практики решения прикладных 


задач. Однако это позволило во всех экспериментах достичь 


асимптотических зон и


установить асимптотическое поведение исследуемых величин для


их сравнения с теоретическими оценками.








Представленный в данной работе алгоритм допускает глубокое


распараллеливание. Основной объем вычислений может производиться


независимо для каждого конечного элемента, каждой стороны и


подзадачи для неизвестных в вершинах. Если для локальных задач


на конечных элементах применяется вторичный МДО, то возможно


дальнейшее распараллеливание. С указанным свойством связано не


менее важное другое: описываемый алгоритм МДО


удобен для использования в адаптивных режимах. Он не


теряет в быстродействии, если порядки пространств внутренних


координатных функций на каждом элементе и пространств следов 


конечно-элементных функций на каждой стороне различны.








\section{Алгоритм МДО и его основные компоненты}





1.1. {\it Иерархическая $p$-версия}. Рассмотрим


обобщенную постановку задачи Дирихле для эллиптического уравнения второго 


порядка в области $\Omega\subset R^2$: 


найти функцию $u\in\W$ такую, что


\begin{equation*}


%%\label{bvp}


a_\Omega(u,v)-(f,v)=0\ \ \text{для любого}\ \ v\in\W.


\end{equation*}


Билинейная форма $a_\Omega(u,v)$ предполагается симметричной,


ограниченной и коэрцитивной. Граница области $\Omega$ состоит из


конечного числа гладких кривых, углы между которыми лежат в


интервале $(0,2\pi)$. Тогда для любого достаточно малого параметра


сетки $h$ область $\ov{\Omega}$ представима в виде


объединения $R={\cal O}(h^{-2})$ непересекающихся совместных


криволинейных четырехугольников $\mathop{\cup}\limits_{r=1}^R\ov{\Pi}_r$,


причем совместные явно заданные отображения $x=x^{(r)}(y)$


квадрата $\overline{\Pi}=[-1,1]\times [-1,1]$ на четырехугольники


$\overline{\Pi}_r$ невырождены и удовлетворяют обобщенным условиям


квазиоднородности (напр., \cite{KJ1999}, \cite{KJ1999a}). Базисный конечный


элемент $\cal E$ есть квадрат $\ov{\Pi}$ с базисом 


${\cal M}=\{ {\cal L}_{i,j}$, $0\leq i,j\leq p\}$, где


${\cal L}_{i,j}(x_1,x_2)={\cal L}_i(x_1) {\cal


L}_j(x_2)$, функции ${\cal L}_j$ для $j>1$ являются


интегрированными полиномами Лежандра с нормировочным множителем


$$ 


{\cal L}_j(x)=\beta_j \int_{-1}^x L_{j-1}(t)dt,\quad


\beta_j=\frac{1}{2}\sqrt{(2j-3)(2j-1)(2j+1)},


$$


а ${\cal L}_0(x)=0{,}5(1+x)$, ${\cal L}_1(x)=0{,}5(1-x)$


--- линейные ``узловые'' функции.





Представим множество ${\cal M}$ в виде объединения подмножеств 


{\em внутренних координатных функций} ${\cal M}_\ri = \{ {\cal


L}_{i,j}$, $2\leq i,j\leq p\}$, {\em координатных функций сторон} 


${\cal M}_\rii= \{ {\cal 


L}_{i,j}$, $i=0,1$, $2\le j\le p$ или $j=0,1$, $2\le i\le p \}$


и {\em координатных функций вершин} 


${\cal M}_\riii=\{ {\cal L}_{i,j}$, $i=0,1$, $j=0,1 \}$. Введем обычным


способом конечноэлементное пространство


$$ 


H^0(\Omega)=\{ u\in C(\ov{\Omega})\mid u(x^{(r)}(y))\in


\Span({\cal M}), \ \ y\in \Pi, \ \ r=1,\dotsc, R


\} \cap \W


$$ 


и его базис $\Phi=\{ \phi_i \}_{i=1,\dotsc, {\cal N}}$ 


разобъем на три подмножества $\Phi_\ri$,


$\Phi_\rii$ и $\Phi_\riii$. Подмножество внутренних базисных 


функций $\Phi_\ri$ состоит из функций $\phi_i$, каждая из которых 


отличается от нуля только на одном элементе $\Pi_r$ и


$\phi_i(x^{(r)}(y)) \in {\cal M}_\ri$. В подмножество базисных функций 


сторон $\Phi_\rii$ входят функции $\phi_i$, которые отличны от нуля на 


паре конечных элементов $\Pi_r$, $\Pi_t$ с общей


стороной и $\phi_i(x^{(r)}(y)), \phi_i(x^{(t)}(y)) \in {\cal M}_\rii$.


Подмножество базисных функций вершин $\Phi_\riii$


составляют функции $\phi_i$, не равные нулю только на конечных


элементах $\Pi_r$ с общей вершиной, на которых 


$\phi_i(x^{(r)}(y)) \in {\cal M}_\riii$. Аналогично разобьем неизвестные 


и обозначим через ${\bold u}_\ri\in \Bbb R^{{\cal N}_\ri}$,


${\bold u}_\rii\in \Bbb R^{{\cal N}_\rii}$,


${\bold u}_\riii\in {\Bbb R}^{{\cal N}_\riii}$ векторы внутренних 


неизвестных и неизвестных сторон и вершин соответственно, где


${\mc N}={\mc N}_\ri+{\cal N}_\rii+{\cal N}_\riii$ --- общее число 


неизвестных. Принятому упорядочению соответствует представление 


конечноэлементной системы алгебраических уравнений 


\begin{equation}


\label{slau}


{\bold K}{\bold u} = {\bold f}


\end{equation}


в блочном виде


\begin{equation*}


%% \label{Kf}


{\bold K}=\left(\begin{array}{lll}


{\bold K}_\ri & {\bold K}_{\ri,\rii} & {\bold K}_{\ri,\riii} \\


{\bold K}_{\rii,\ri} & {\bold K}_{\rii} & {\bold K}_{\rii,\riii} \\


{\bold K}_{\riii,\ri} & {\bold K}_{\riii,\rii} & {\bold K}_{\riii}


\end{array}


\right), \qquad {\bold u}=\left(\begin{array}{l}


{\bold u}_\ri \\ {\bold u}_{\rii} \\ {\bold u}_{\riii}


\end{array}


\right), \qquad {\bold f}=\left(\begin{array}{l}


{\bold f}_\ri \\ {\bold f}_{\rii} \\ {\bold f}_{\riii}


\end{array}


\right).


\end{equation*}








1.2. {\it Структура предобусловливателя МДО}. \label{MDO} Следуя


\cite{KJ1999}, \cite{KJ1999a} для построения предобусловливателя МДО введем 


более простую матрицу, чем ${\bold K}$. Свяжем с каждым элементом 


$\Pi_r$ матрицу ${\bold A}_r\equiv {\bold A} = \wh{{\bold A}}_1


+ h^2 \wh{{\bold A}}_0$ или ${\bold A}_r\equiv{\bold A}=\wh{{\bold A}}_1$


где $ \wh{{\bold A}}_1$ --- матрица жесткости базисного элемента ${\cal E}$,


индуцируемая билинейной формой 


$\wh{a}_\Pi(w,v)=\int\limits_\Pi(\triangledown w\cdot\triangledown v +


wv)dx$, a $ \wh{{\bold A}}_0$ --- его матрица массы при локальной 


нумерации неизвестных.


Обозначим через ${\boldsymbol \Lambda}$ глобальную матрицу, получающуюся в


результате процедуры сборки (аналогичной сборке $\bold K$) матриц


${\bold A}_r$, и представим ее в блочном виде с учетом разделения


неизвестных на внутренние, неизвестные сторон и вершин


\begin{equation}


\label{Lambda}


{\bi \Lambda}=\left(\begin{array}{lll}


{\bi \Lambda}_\ri & {\bi \Lambda}_{\ri,\rii} & {\bi \Lambda}_{\ri,\riii} \\


{\bi \Lambda}_{\rii,\ri} & {\bi \Lambda}_{\rii} & {\bi \Lambda}_{\rii,\riii} \\


{\bi \Lambda}_{\riii,\ri} & {\bi \Lambda}_{\riii,\rii} & {\bi \Lambda}_{\riii}


\end{array}


\right).


\end{equation}


Спектральная эквивалентность матриц $\bi \Lambda$ и ${\bold K}$ 


следует~\cite{K1977} из условий обобщенной квазиоднородности. 


Кроме того, матрица


$$


{\bi \Lambda}^{(0)}=\left(\begin{array}{lll}


{\bi \Lambda}_\ri & {\bi \Lambda}_{\ri,\rii} & \bold O \\


{\bi \Lambda}_{\ri,\ri} & {\bi \Lambda}_{\ri} & \bold O \\


\bold O & \bold O & {\bi \Lambda}_{\riii}


\end{array}


\right)=\left(\begin{array}{lll}


\bi \Lambda^{(\ri)} & \bold O \\


\bold O & {\bi \Lambda}_{\riii}


\end{array}


\right),


$$


полученная при помощи обнуления блоков взаимодействия


функций вершин с остальными, является \cite{KJ1999}, \cite{KJ1999a} почти


спектрально-эквивалентной ${\bold K}$:


$$


(1+\ln p)^{-1} {\bi \Lambda}^{(0)} \prec {\bold K} \prec


{\bi \Lambda}^{(0)}.


$$


Знак $\prec$ используем для записи неравенств, выполняющихся с


точностью до не зависящих от $h$ и $p$ постоянных.





Глобальный МДО-предобусловливатель $\bc K$ для


${\bi\Lambda}^{(0)}$ определяем через обратный к нему, приняв


\begin{equation} \label{globprec}


\bc K^{-1}=\bi {\cal K}_\ri^{(+)} +


\bc K_{\rii}^{(+)} + \bc K_{\riii}^{(+)}.


\end{equation}


Основной версии пакета соответствуют


\begin{equation*}


%% \label{plzd}


{\bc K}_\ri^{(+)}= \left( \begin{array}{lll}


{\bc K}_\ri^{-1} &\bold O & \bold O \\


\bold O & \bold O & \bold O \\


\bold O & \bold O & \bold O 


\end{array}


\right), \quad \bc K_{\riii}^{(+)}= \left(\begin{array}{lll}


\bold O & \bold O & \bold O \\


\bold O & \bold O & \bold O \\


\bold O & \bold O & {\bc K }_{\riii}^{-1} 


\end{array}


\right),


\end{equation*}


где ${\bc K }_{L}^{-1}$ --- неявно заданные матрицы,


являющиеся результатом применения неточных итерационных солверов


с фиксированным и небольшим числом итераций


для решения систем с матрицами ${\bi \Lambda}_L$. 


В ${\bc K}_\ri^{(+)}$ и ${\bc K}_{\riii}^{(+)}$


можно, вообще говоря, использовать ${\bc K}_L^{-1}={\bi \Lambda}_L^{-1}$,


$L=\ri,\riii$, т.\,е. псевдообратные матрицы к соответствующим блокам 


матрицы ${\bi \Lambda}$, однако это предполагает


точные методы решения систем с матрицами ${\bi \Lambda}_L$. Пусть ${\bold


A}_\ri$ --- блок матрицы ${\bold A}$, определенный на внутренних 


неизвестных. Поскольку


${\bi \Lambda }_{\ri}=\diag \{ {\bold A}_\ri,


{\bold A}_\ri, \dotsc, {\bold A}_\ri \}$, то решение системы с


матрицей ${\bi \Lambda}_\ri$ требует решения $R$ систем с матрицей 


${\bold A}_\ri$ и соответственно предобусловливатель локальных задач 


Дирихле определяется как блочно-диагональная матрица с $R$ 


одинаковыми блоками ${\bc K }_{\ri}^{-1}=\diag \{


\wt {{\bi {\mc A}}}{}_\ri^{-1},


\wt{{\bi {\mc A}}}{}_\ri^{-1}, \dotsc, \wt{{\bi {\mc A}}}{}_\ri^{-1}\}$. Здесь


$\wt{{\bi {\mc A}}}{}_\ri^{-1}$ --- оператор, являющийся результатом 


применения некоторого вторичного итерационного процесса с небольшим 


числом итераций для решения системы с матрицей ${\bi {\mc A}}_\ri$. 


Всюду в статье используем волну в обозначении операторов, 


порожденных вторичным


итерационным процессом для решения системы линейных уравнений.





Определение предобусловливателя интерфейсных задач $\bi{\cal


K}_{\rii}^{(+)}$ связано с разложением


\begin{equation*}


%%\label{LambdaI}


\bi \Lambda^{(\ri)}=\left(\begin{array}{ll}


\bold I & \bold O \\


\bi \Lambda_{\rii,\ri}\bi \Lambda_\ri^{-1} & \bold I


\end{array}\right)\,


\left(\begin{array}{ll}


\bi \Lambda_\ri & \bold O \\


\bold O & \bold S


\end{array}\right)\,


\left(\begin{array}{ll}


\bold I & \bi\Lambda_\ri^{-1}\bi\Lambda_{\ri,\rii} \\


\bold O & \bold I \end{array}\right),


\end{equation*}


где ${\bold S}={\bi \Lambda}_{\ri} - {\bi \Lambda}_{\rii,\ri}


{\bi\Lambda}_\ri^{-1}{\bi\Lambda}_{\ri,\rii}$ ---


дополнение Шура. В дальнейшем будем использовать блочно-диагональный 


предобусловливатель $\bi {\cal S}$ для дополнения Шура с одинаковыми 


блоками для каждой из сторон триангуляции (при условии однотипной 


нумерации неизвестных каждой из сторон), т.\,е. ${\bc S}=\diag \{ {\bc


S}_0, {\bc S}_0, \dotsc, {\bc S}_0\}$. При этом он удовлетворяет 


неравенствам


\begin{equation}\label{specshur}


(1+\ln p)^{-1} \bi{\cal S} \prec {\bold S}


\prec (1+\ln p) \: \bi{\cal S}.


\end{equation}





Пусть $\bi {\cal P} : {\Bbb R}^{{\cal N}_{\rii}}


\rightarrow {\Bbb R}^{{\cal N}_{\ri} + {\cal N}_{\rii}}$


--- оператор продолжения такой, что для любого вектора ${\bold


v}_{\rii}\in {\Bbb R}^{{\cal N}_{\rii}}$ выполняется


неравенство


\begin{equation}\label{prolongoper}


\| \bi{\cal P} {\bold v}_{\rii} \|_{{\bi 


\Lambda}^{(\ri)}}\le\gamma_{\cal P}\|{\bold v}_{\rii} \|_{\bold S},


\end{equation}


c $\gamma_{\cal P}$ постоянной или слабо зависящей от $p$. 


Полагаем $ \bi {\cal


K}_{\rii}^{(+)} = \bi{\cal P} \bi {\cal S}^{-1} \bi {\cal P}^\top$.


Обозначим через ${\bc P}_r$ ограничение оператора ${\bc P}$ на


элемент с номером $r$. В силу свойств матриц ${\bi \Lambda}_\ri$ и ${\bc


S}$ все эти ограничения одинаковы при условии идентичности


локальных нумераций внутренних координатных функций на всех


элементах и координатных функций сторон на всех сторонах,


т.\,е. ${\bc P}_r\equiv {\bc P}_{_{


E}}:\Bbb {R}^{4(p-1)}\rightarrow \Bbb{R}^{(p-1)(p-1)+4(p-1)}$ --- 


оператор продолжения для базисного элемента. Более того,


если для блоков матрицы ${\bold A}$ ввести аналогичную 


(\ref{Lambda}) индексацию и использовать обозначения


\begin{equation*}%%%\label{}


{\bold A}^{(\ri)} = \begin{pmatrix}


{\bold A}_\ri & {\bold A}_{\ri,\rii} \\


{\bold A}_{\rii,\ri} & {\bold A}_{\rii}


\end{pmatrix}, \quad


\Bbb S_E= {\bold A}_{\rii}-{\bold A}_{\rii,\ri}{\bold


A}_\ri^{-1}{\bold A}_{\ri,\rii},


\end{equation*}


то (\ref{prolongoper}) эквивалентно неравенству


\begin{equation}\label{prolongoperE}


\| {\bc P}_E {\bold v}_{\rii} \|_{{\bold


A}^{(\ri)}} \le \gamma_{\cal P} \| {\bold v}_{\rii} \|_{\Bbb S_E}.


\end{equation}





Описанный способ построения МДО-предобусловливателей


использовался в \cite{KI1996}--\cite{KJ1999a}.


Справедлива теорема (\cite{KJ1999a}, теорема~5.2),


свидетельствующая о том, что предобусловливатель ${\bc K}$


почти оптимален в смысле относительного числа обусловленности.





\begin{thms}\label{TeorGlobPrec}


Пусть $(1+\ln p)^{-1}{\bc K}_L \prec


\Lambda_L \prec {\bc K}_L $ для $L=\ri,\riii$ и выполняются


неравенства $(\ref{specshur})$ и неравенство $(\ref{prolongoperE})$ с


постоянной $\gamma_{\cal P}$.  Тогда


$$


(1+\ln p)^{-2}\bi{\cal K} \prec {\bold K} \prec (1+\ln p) \bi{\cal K}.


$$


\end{thms}





Для решения системы (\ref{slau}) применим 


ПМСГ с предобусловливателем, обратный к которому


имеет вид~(\ref{globprec}),





${\bold u}^{(1)} := 0$, ${\bold r}^{(1)} := {\bold f}$,


${\bold p}^{(0)} := 0$, ${\bold r}^{(0)} := 0$,


${\beta}_1 := 0$





\underline{{\bf for}} $k:=1,2,\dotsc$





$\quad {\bold z}^{(k)}:=\bi{\cal K}^{-1} {\bold r}^{(k)}$





\quad \underline{{\bf if}} $k>1$





$\qquad \beta_k:= [({\bold z}^{(k)})^\top ( {\bold r}^{(k)})]/


[({\bold z}^{(k-1)})^\top{\bold r}^{(k-1)}]$





\quad \underline{{\bf endif}}





\quad ${\bold p}^{(k)} := {\bold z}^{(k)} + \beta_k {\bold p}^{(k-1)}$,


$\alpha_k :={\displaystyle


\frac{({\bold z}^{(k)})^\top {\bold r}^{(k)}}{({\bold p}^{(k)})^\top


{\bold K}{\bold p}^{(k)}}}$,


${\bold r}^{(k+1)} := {\bold r}^{(k)}-\alpha_k {\bold K}{\bold p}^{(k)}$,


${\bold u}^{(k+1)} := {\bold u}^{(k)}+\alpha_k {\bold p}^{(k)}$





\underline{{\bf endfor}}.





В последнее время проблеме устойчивости итерационных методов решения систем


линейных уравнений при предобусловливании посредством вторичных


точных и неточных солверов уделяется много внимания. В ряде работ


она рассматривалась с чисто алгебраической точки зрения


(напр., \cite{GGK}--\cite{GO}). В ~\cite{GE} был предложен вариант 


ПМСГ, менее чувствительный к


ошибкам решения системы с предобусловливателем. В этом варианте,


который также реализован в пакете программ и использовался в наших


численных экспериментах, итерационный параметр $\beta_k$


вычисляется по формуле


\begin{equation}\label{betak}


\beta_k:=({\bold z}^{(k)})^\top ( {\bold r}^{(k)}-{\bold r}^{(k-1)}) /


[({\bold z}^{(k-1)})^\top{\bold r}^{(k-1)}]\:.


\end{equation}





Далее рассмотрим основные компоненты МДО-предобусловливателя


$\bi{\cal K}$, которые определим таким образом, что


вычислительная стоимость решения системы с матрицей $\bi{\cal


K}$ будет почти оптимальной.


\medskip





1.3. {\it Предобусловливание локальных задач Дирихле}. \label{ss:1.3}


Как следует из предыдущей секции, умножение ${\bc K}_\ri^{(+)}{\bold v}_\ri$


требует лишь решения $R$ систем с матрицей ${\bi {\mc A}}_\ri$.


Известные приемы построения быстрых предобусловливателей-солверов


${\bi {\mc A}}_\ri$


(см. \cite{B2002}, \cite{BSS2002}, \cite{K2001}--\cite{Korneev:02b})


основаны на спектрально эквивалентном предобусловливателе 


конечно-разностного вида, полученном в \cite{KI1996},


который будем называть {\em базовым}.


Для простоты изложения считаем $p=2N+1$ с целым $N$ и разобьем множество 


внутренних координатных функций базисного элемента на четыре


подмножества в соответствии с четностью-нечетностью их порядков по каждой 


из двух переменных: ${\cal M}_\ri={\cal M}^{ee}_\ri\cup


{\cal M}^{eo}_\ri\cup {\cal M}^{oe}_\ri\cup {\cal M}^{oo}_\ri$,


где


\begin{align*}


{\cal M}^{ee}_\ri& =\{ {\cal L}_{2,2}, {\cal L}_{2,4},


\dotsc,{\cal L}_{2,2N},


{\cal L}_{4,2}, {\cal L}_{4,4},\dotsc,{\cal L}_{4,2N},\dotsc,


{\cal L}_{2N,2}, {\cal L}_{2N,4},\dotsc,{\cal L}_{2N,2N}\},\\


%


{\cal M}^{eo}_\ri& =\{ {\cal L}_{2,3}, {\cal L}_{2,5},\dotsc,


{\cal L}_{2,2N+1},


{\cal L}_{4,3}, {\cal L}_{4,5},\dotsc,{\cal L}_{4,2N+1},\dotsc,


{\cal L}_{2N,3}, {\cal L}_{2N,5},\dotsc,{\cal L}_{2N,2N+1}\},\\


%


{\cal M}^{oe}_\ri& =\{ {\cal L}_{3,2}, {\cal L}_{3,4},\dotsc,


{\cal L}_{3,2N},


{\cal L}_{5,2}, {\cal L}_{5,4},\dotsc,{\cal L}_{5,2N},\dotsc,


{\cal L}_{2N+1,2}, {\cal L}_{2N+1,4},\dotsc,{\cal L}_{2N+1,2N}\},\\


%


{\cal M}^{oo}_\ri& =\{ {\cal L}_{3,3}, {\cal L}_{3,5},\dotsc,


{\cal L}_{3,2N+1},{\cal L}_{5,3}, {\cal L}_{5,5},\dotsc,


{\cal L}_{5,2N+1},\dotsc, {\cal L}_{2N+1,3},


{\cal L}_{2N+1,5},\dotsc,{\cal L}_{2N+1,2N+1}\}.


\end{align*}


После этого матрица жесткости внутренних функций ${\bold A}_\ri$


базисного элемента становится блочно-диагональной: ${\bold A}_\ri


= \diag\{{\bold A}_{e,e},{\bold A}_{e,o},{\bold A}_{o,e},


{\bold A}_{o,o}\}$ (см. \cite{KI1996}). Блоки базового блочно-диагонального


предобусловливателя


${\bc A}_\ri = \diag\{ {\bc A}_{e,e},{\bc A}_{e,o},{\bc


A}_{o,e},{\bc A}_{o,o} \}$ определим соотношениями


\begin{equation*}


%%\label{Gab}


\bi{\cal A}_{a,b}=\bi{\cal


D}_a\otimes(\bi\Delta+\bi{\cal D}_b^{-1}) +


(\bi\Delta+\bi{\cal D}_a^{-1})\otimes\bi{\cal D}_b,


\end{equation*}


в которых $a$, $b$ принимают значения $e$, $o$ и


\begin{gather*}


\bi{\cal D}_e=\diag\{4i^2\}_{i=1,2,\dotsc,N},\quad \bi{\cal


D}_o=\diag\{4i^2+4i+1\}_{i=1,2,\dotsc,N}, \\


%


\bi\Delta=\frac12


\begin{pmatrix}


2 & -1 & & & & & \\


-1 & 2 & -1& & & 0 & \\


& & \dotsc & \dotsc & & & \\


& 0 & & & -1 & 2 &-1 \\


& & & & & -1 &2 


\end{pmatrix}.


\end{gather*}


Этот предобусловливатель спектрально-эквивалентен матрице жесткости 


внутренних функций базисного элемента \cite{KJ1997}: 


${\bc A}_\ri\prec {\bold A}_\ri \prec {\bc A}_\ri$.


Кроме того, все блоки $\bi{\cal A}_{a,b}$ спектрально-эквивалентны друг


другу, поэтому можно также положить $\bi{\cal A}=\diag\{\bi{\cal


A}_{a,a},\bi{\cal A}_{a,a}, \bi{\cal A}_{a,a}, \bi{\cal A}_{a,a}\}$


c $a=e$ или $a=o$. Таким образом, для получения быстрого солвера для 


матрицы $\bi{\cal A}_\ri$ достаточно получить


такой солвер для $\bi{\cal A}_{e,e}$.





В \cite{KI1996}, как и в \cite{IvanovKorneev:95}, ставилась цель получить 


предобусловливатель конечно-разностного вида. Очевидно, 


${\bi{\cal A}}_{e,e}$ является матрицей конечно-разностной аппроксимации 


на квадратной сетке с шагом $\hslash=1/N$ (и пятиточечным


шаблоном) эллиптического уравнения второго порядка


$$


L_{{\bi{\cal A}}}u\equiv-2(\xi_1^2 u_{\xi_2\xi_2}''+\xi_2^2 u_{\xi_1\xi_1}'')+


(\xi_1^2/\xi_2^2 + \xi_2^2/\xi_1^2)u =


F(\xi), \quad \xi=(\xi_1,\xi_2)\in\pi_1, \quad u|_{\partial \pi_1}=0,


$$


в квадрате $\pi_1:=(0,1)\times(0,1)$. Коэффициенты этого уравнения 


вырождаются, а коэффициенты при младших


членах имеют, кроме того, особенности. Это явилось причиной того,


что быстрые солверы для матрицы ${\bi {\mc A}}_\ri$ появились лишь 


в последние годы в статьях \cite{B2002}, \cite{K2001}--\cite{Korneev:02b}


и препринте \cite{BSS2002}. Они относятся к трем типам:


МДО для конечно-разностной аппроксимации, многосеточный метод и метод,


основанный на специальной многоуровневой вэйвлетной декомпозиции 


пространства кусочно-линейных функций на равномерной сетке. Во всех 


из них сначала строятся


близкие к ${\bi {\mc A}}_\ri$ модифицированные предобусловливатели, более


приспособленные к конкретному солверу. Например, для построения МДО


предобусловливателя-солвера в \cite{K2001}, \cite{K2002} была введена


более крупная ортогональная декомпозиционная сетка, сгущающаяся по


каждой из переменных $\xi_1$, $\xi_2$ специальным образом к нулю


так, что на каждой ее ячейке коэффициенты отличаются от постоянных


на одну и ту же (относительно) часть, т.\,е. фактически ${\bi


{\mc A}}_\ri$ заменялся спектрально эквивалентным


предобусловливателем ${\bold B}_\ri$, соответствующим


дифференциальному оператору $L_{\bold B}$ с кусочно-постоянными


коэффициентами, для которого строился МДО. В \cite{B2002}, \cite{BSS2002}


многоуровневые алгоритмы строились для упрощенного с потерей спектральной


эквивалентности предобусловливателя


\begin{equation}\label{Gbeu}


{\Bbb A}_\ri = \diag\{{\Bbb A}_{e,e},{\Bbb A}_{e,o},


{\Bbb A}_{o,e},{\Bbb A}_{o,o}\}, \quad {\Bbb A}_{a,b}=\bi{\cal


D}_a\otimes \bi\Delta + \bi\Delta \otimes\bi{\cal D}_b,


\end{equation} 


для которого $\cond({\Bbb A}_\ri^{-1}{\bold A}_\ri) \prec (1+ \log p)$. 


Он является конечно-разностной аппроксимацией оператора краевой задачи


$$


L_{\Bbb A} u \equiv -2(\xi_1^2 u_{\xi_2\xi_2}''+\xi_2^2 u_{\xi_1\xi_1}'')=


F(\xi), \quad\xi\in\pi_1,\quad u|_{\partial \pi_1}=0.


$$


Один из модулей пакета реализует также обобщение многоуровневого алгоритма на


случай произвольного $p$.\footnote{Этот алгоритм и модуль разработаны


А.Б. Смирновым и Е.Н. Смирновой.}


\setcounter{footnote}{0}





Общая вычислительная стоимость предобусловливателей-солверов


${\Bbb A}_\ri$ и ${\bold B}_\ri$ для матрицы ${\bi {\mc A}}_\ri$


колеблется в зависимости от версии солвера от $p^2(1+\ln p)^{0{,}5}$


до $p^2(1+\ln p)^{4{,}5}$. За более подробным описанием и результатами 


исследования этих солверов мы отсылаем


к цитированным выше работам.


\medskip





1.4. {\it Предобусловливание дополнения Шура и продолжения с


интерфейса}. \label{PSPR}


Для построения $\bi{\cal K}_{\rii}^{(+)}$ нужны: предобусловливатель


$\bi{\cal S}$ для дополнения Шура ${\bold S}$


и оператор продолжения $\bi{\cal P}$,


удовлетворяющие условиям (\ref{specshur}) и (\ref{prolongoperE}).


Для определения блока


$\bi{\cal S}_0$ матрицы $\bi{\cal S}$ достаточно, следуя


\cite{KJ1997}, рассмотреть предобусловливание дополнения 


Шура $\Bbb{ S}_{E}$ базисного элемента. Будем считать, 


не нарушая общности, что $p=2N+1$. Переупорядочим координатные функции 


сторон базисного элемента, представив множество функций сторон в виде


${\cal M}_{\rii}={\cal M}^{(1)}_{\rii}\cup {\cal M}^{(2)}_{\rii}


\cup {\cal M}^{(3)}_{\rii}\cup {\cal M}^{(4)}_{\rii}$ так, что каждое 


из ${\cal M}^{(s)}_{\rii}$ будет содержать координатные функции $s$-ой 


стороны, пронумерованные с учетом четности-нечетности индексов,


\begin{equation} \label{sideren}


\begin{split}


{\cal M}^{(1)}_{\rii}&={\cal M}^{(1,e)}_{\rii}


\cup{\cal M}^{(1,o)}_{\rii},\ \ {\cal


M}^{(1,e)}_{\rii}=\{ {\cal L}_{2,1}, {\cal L}_{4,1}, ...,


{\cal L}_{2N,1}\},\ \ {\cal M}^{(1,o)}_{\rii}=


\{ {\cal L}_{3,1}, {\cal L}_{5,1}, ...,


{\cal L}_{2N+1,1}\},\\


%


{\cal M}^{(2)}_{\rii}&={\cal M}^{(2,e)}_{\rii}


\cup{\cal M}^{(2,o)}_{\rii},\ \ {\cal


M}^{(2,e)}_{\rii}=\{ {\cal L}_{1,2}, {\cal L}_{1,4}, ...,


{\cal L}_{1,2N}\},\ \ {\cal M}^{(2,o)}_{\rii}=


\{ {\cal L}_{1,3}, {\cal L}_{1,5}, ...,{\cal L}_{1,2N+1}\},\\


%


{\cal M}^{(3)}_{\rii}&={\cal M}^{(3,e)}_{\rii}\cup


{\cal M}^{(3,o)}_{\rii},\ \ {\cal M}^{(3,e)}_{\rii}=


\{ {\cal L}_{2,0}, {\cal L}_{4,0}, ..., {\cal L}_{2N,0}\}, 


\ \ {\cal M}^{(3,o)}_{\rii}=\{ {\cal L}_{3,0}, {\cal L}_{5,0}, ...,


{\cal L}_{2N+1,0}\},\\


%


{\cal M}^{(4)}_{\rii}&={\cal M}^{(4,e)}_{\rii}


\cup{\cal M}^{(4,o)}_{\rii},\ \ {\cal M}^{(4,e)}_{\rii}=


\{ {\cal L}_{0,2}, {\cal L}_{0,4}, ..., {\cal L}_{0,2N}\}, 


\ \ {\cal M}^{(4,o)}_{\rii}=\{ {\cal L}_{0,3}, {\cal L}_{0,5}, ...,


{\cal L}_{0,2N+1}\}.


\end{split}


\end{equation}


Предобусловливатель для $\Bbb{S}_E$ выбираем в виде 


$\bi{\cal S}_E=\diag[ \bi{\cal S}_0, \bi{\cal S}_0,


\bi{\cal S}_0, \bi{\cal S}_0 ]$, где каждый блок


$\bi{\cal S}_0$ отвечает одной из сторон базисного элемента.


Определение $\bi{\cal S}_0$ связано с различными


представлениями полинома $q(x)$ степени $2N+1$ на отрезке


$$


q(x)=\sum_{i=0}^{2N+1} a_i {\cal L}_i(x),\quad q(x)=


\sum_{i=0}^{2N+1} b_i x^i, \quad


q(\cos\phi)=\sum_{i=0}^{2N+1} c_i \cos i\phi.


$$


Введем матрицу $\bi{\cal D}=\diag\{0,1,\dotsc, 2N+1\}$ 


и верхние треугольные матрицы


$\bi{\cal B}$, $\bi{\cal C} $, которые соответствуют преобразованиям


$$


\bi{\cal B}{\bold a}={\bold b},\ \ \bi{\cal C}{\bold b}=


{\bold c},\ \ {\bold a}=\{a_i\}_{i=0,1,\dotsc,2N+1},


\ \ {\bold b}=\{b_i\}_{i=0,1,\dotsc,2N+1}, \ \ {\bold


c}=\{c_i\}_{i=0,1,\dotsc,2N+1}.


$$


Пусть $\bi{\cal W}=\bi{\cal C}\bi{\cal B}$. Матрица $\bi{\cal S}_0$ 


определяется как матрица $\ln(1+p)\bi{\cal W}^\top \bi{\cal D} \bi{\cal W}$


с вычеркнутыми строками и столбцами, соответствующими $a_0$ и $a_1$,


т.\,е. первым двум функциям ${\cal L}_0$ и ${\cal L}_1$.





\begin{thms}[{\series{m}\shape{n}\selectfont \cite{KJ1999}, теорема 3.1}]


\label{TeorSE}


Справедливы неравенства $(1{+}\ln p)^{-1}\bi{\cal S}_ E {\prec}


{\Bbb{S}}_E{\prec}(1{+}\ln p)^{k-1}


\bi{\cal S}_ E$ при ${\bold A}=\wh{{\bold A}}_1$, $k=3$


и при ${\bold A}=\wh{{\bold A}}_1 + h^2\wh{{\bold A}}_0$, $k=2$.


\end{thms}





%%%% ^^^^^ remove extra { }








Заметим, что из теоремы \ref{TeorSE} следует (\ref{specshur}),


если при сборке $\bi \Lambda$ использовались матрицы \linebreak


${\bold A}_r\equiv \wh{{\bold A}}_1$.





Остановимся на некоторых особенностях реализации предобусловливателя 


$\bi{\cal S}$. Принятая


перенумерация координатных функций сторон (\ref{sideren}) приводит


к тому, что матрицы $\bi{\cal W}$ и $\bi{\cal D}$


распадаются на два независимых подблока для неизвестных,


соответствующих четным и нечетным базисным функциям. Cледовательно,


\begin{equation*}


%%\label{S0}


\bi{\cal W}^\top \bi{\cal D} \bi{\cal W}=


\begin{pmatrix}


\bi{\cal W}^{(e)} & \bold O \\


\bold O & \bi{\cal W}^{(o)}


\end{pmatrix}


^\top


\begin{pmatrix}


\bi{\cal D}^{(e)} & \bold O \\


\bold O & \bi {\cal D}^{(o)}


\end{pmatrix}


\begin{pmatrix}


\bi{\cal W}^{(e)} & \bold O \\


\bold O & \bi{\cal W}^{(o)}


\end{pmatrix}


,


\end{equation*}


где $\bi{\cal D}^{(e)}=\diag\{0,2,\dotsc, 2N\}$,


$\bi{\cal D}^{(o)}=\diag\{1,3,\dotsc, 2N+1\}$. 


Таким образом, предобусловливание интерфейсной задачи на каждом шаге


внешнего итерационного процесса требует решения $2{\cal N}_s$ систем 


вида $\bi{\cal F}{\bold y}={\bold d}$, где $\bi{\cal F}$ получаем 


вычеркиванием первой строки из матрицы


\begin{equation*}


%%\label{SF}


\ov{\bc S_0} = \ov{\bc W}^\top\ov{\bc D}


\ov{\bc W},


\end{equation*}


в которой либо $\ov{\bi{\cal W}}=\bi{\cal W}^{(e)}$ и


$\ov{\bi{\cal D}}=\bi{\cal D}^{(e)}$, либо $\ov{\bi{\cal W}}=\bi{\cal


W}^{(o)}$ и $\ov{\bi{\cal D}}=\bi{\cal D}^{(o)}$. 


Простой прием помогает избежать перемножения матриц с последующим 


вычеркиванием первой строки и столбца и решением систем с


полностью заполненными матрицами. Обозначим через 


$\bold z$ и $\bold u$ решения вспомогательных систем


\begin{equation}


\label{helpsys}


\ov{\bi{\cal W}}^\top\ov{\bi{\cal D}}\,


\ov{\bi{\cal W}}{\bold


z}=\{\,\underbrace{1,0,\dotsc,0}_{N+1}\,\}^\top,\quad


\ov{\bi{\cal W}}^\top\ov{\bi{\cal D}}\,


\ov{\bi{\cal W}}{\bold u}=


\begin{bmatrix}


0 \\ {\bold d} 


\end{bmatrix}.


\end{equation}


Первую из них достаточно решить один раз перед началом итерационного


процесса. Компоненты решения системы $\bi{\cal F}{\bold y}={\bold d}$ 


находятся по формулам $y_k=u_{k+1}-(u_1/z_1) z_{k+1}$ для 


$k=1,2,\dotsc, N$. Следовательно, обращение предобусловливателя 


интерфейсных задач $\bi{\cal S}$ сводится к решению систем с треугольными и


диагональной матрицами для каждой стороны триангуляции, и его вычислительная 


стоимость составляет


${\cal N}_s{\cal O}(p^2)={\mc O}(Rp^2)$, где ${\cal N}_s$ --- число сторон.





Выбор оператора продолжения $\bi{\cal P}$, эквивалентно $\bi{\cal P}_E$, 


существенно определяет эффективность интерфейсной компоненты


$\bi{\cal K}_{\rii}^{(+)}$. Положим 


$\bi{\cal P}^\top_E=(-{\bold A}_{\rii,\ri}\wt{\bold A}


_{\ri,\mathrm{pr}}^{-1}\mid {\bold I}_{\rii})$,


где $\wt{\bold A} _{\ri,\mathrm{pr}}^{-1}$ является оператором, 


порожденным итерационным процессом


\begin{equation}\label{inexact1}


{\bold w}^{k+1}={\bold w}^{k}+ \sigma_k {\bc A}_1^{-1}


({\bold A}_1{\bold w}^{k}-


{\bold v}_{\ri,r}), \quad {\bold w}^0={\bold 0},


\end{equation}


для приближенного решения систем 


${\bold A}_\ri{\bold w}={\bold v}_{\ri,r}$ для каждого конечного элемента


$r$, в котором $k=0,1,\dotsc,n$ и $\sigma_k $ --- чебышевские 


итерационные параметры, ${\bi {\mc A}}_\ri$ --- быстрый


предобусловлива\-тель-сол\-вер для матрицы ${\bold A}_\ri$.








Таким образом, формально определяется матрица


\begin{equation*}


%\label{inexact2}


\wt{\bold A}_{\ri,\mathrm{pr}}^{-1}=\bigg[{\bold I}-


\prod_{k=1}^{n}({\bold I}-\sigma_k {\bc A }_\ri^{-1} {\bold A}_\ri) \bigg]


{\bold A}_\ri^{-1}.


\end{equation*}


В \cite{KXA2002}, \cite{KXA2002P} доказано, что для выполнения 


свойства~(\ref{prolongoper}) оператора продолжения достаточно сделать 


$n=n_{\mathrm{pr}}(p)={\mc O}((\ln (1+p))/\rho)$ итераций


в~(\ref{inexact1}), где $\rho$ --- коэффициент сжатия невязки на одной 


итерации. Для метода простых итераций


с постоянным итерационным параметром $\sigma_k\equiv \sigma_0={\mc


O}(1)$ имеем


$\rho=(\lambda_{\max}-\lambda_{\min})/(\lambda_{\max}+\lambda_{\min})={\mc


O}(1)$, где $\lambda_{\max}$ и $\lambda_{\min}$ --- максимальное и


минимальное собственные числа обобщенной проблемы ${\bold A}_\ri \bold x


= \lambda {\bi{\cal A}}_\ri \bold x$. Такой выбор оператора


продолжения приводит к тому, что


$$


\bi{\cal K}_{\rii}^{(+)} =


\begin{pmatrix}


-\wt{\bi\Lambda}{}_{\ri,\mathrm{pr}}^{-1}\bi\Lambda_{\ri,\rii} \\


{\bold I}_{\rii}


\end{pmatrix}


\bi{\cal S}^{-1} (-\bi\Lambda_{\ri,\ri}\wt{\bi\Lambda}


{}^{-1}_{\ri,\mathrm{pr}} \mid 


{\bold I}_{\rii}),\quad \wt{\bi\Lambda}{}_{\ri,\mathrm{pr}}^{-1}=


\diag\{\wt{\bold A}{}


_{\ri,\mathrm{pr}}^{-1},\dotsc,\wt{\bold A}{}_{\ri,\mathrm{pr}}^{-1}\}.


$$





Умножение $\bi{\cal K}_{\rii}^{(+)}$ во внешнем итерационном процессе на 


текущем шаге на вектор невязки


$$


{\bold r}= \left(\begin{array}{l}


{\bold r}_\ri \\ {\bold r}_{\rii} \\ {\bold r}_{\riii}


\end{array}


\right) =\left(\begin{array}{l}


{\bold r}^{(\ri)} \\ {\bold r}_{\riii}


\end{array}


\right)


$$


не требует предварительной сборки матриц и


состоит из трех этапов.





1. {\it Сужение}


\begin{equation}\label{restr}


{\bold


w}_{\ri}\leftarrow\wt{\bi\Lambda}{}_{\ri,\mathrm{pr}}^{-1} {\bold


r}_\ri,\quad {\bold z}_{\rii}\leftarrow -\bi\Lambda_{\rii,\ri}{\bold w}_{\ri}


+ {\bold r}_{\rii}.


\end{equation}


Операция $\wt{\bi\Lambda}{}_{\ri,\mathrm{pr}}^{-1} {\bold r}_\ri$ требует 


решения $R$ систем ${\bold A}_\ri {\bold w}_{\ri}^{(r)}={\bold r}_{\ri}^{(r)}$


при помощи $n_{\mathrm{pr}}$ итераций предобусловленного


итерационного процесса (\ref{inexact1}). Умножение 


$\bi\Lambda_{\rii,\ri}{\bold w}_{\ri}$ также выполняется


поэлементно ${\bold A}_{\rii,\ri}{\bold w}_{\ri}^{(r)}$ 


для всех $r$, полученные векторы затем собираются.





2. {\it Предобусловливание дополнения Шура}


$$


{\bold z}_{\rii}\leftarrow \bi{\cal S}^{-1} {\bold w}_{\rii}


$$


требует решения систем вида~(\ref{helpsys}) независимо для каждой


стороны триангуляции.





3. {\it Продолжение}


\begin{equation}


\label{prolong}


\begin{pmatrix}


{\bold v}_\ri \\ {\bold v}_{\rii}


\end{pmatrix}


\leftarrow 


\begin{pmatrix}


\widetilde{\bi\Lambda}{}_{\ri,\mathrm{pr}}^{-1}\bi\Lambda_{\ri,\rii} \\


{\bold I}_{\rii}


\end{pmatrix}


{\bold z}_{\rii}


\end{equation}


заключается в поэлементных умножениях 


$\ov{\bold z}_r\leftarrow{\bold A}_{\ri,\rii}{\bold z}^{(r)}_{\rii}$ с


подвекторами сторон элементов и поэлементном решении систем с матрицей


${\bold A}_\ri {\bold v}_{\ri}^{(r)}=\ov {\bold z}_{\ri}^{(r)}$


так же, как и в п.\,1.





Матрица взаимодействия внутренних функций с координатными функциями сторон 


${\bold A}_{\ri,\rii}$ имеет сильно разреженную структуру


\begin{equation*}


%%\label{tempAI-II}


{\bold A}_{\ri,\rii}=


\begin{pmatrix}


{\bold A}_{ee,\rii} \\ {\bold A }_{eo,\rii} \\


{\bold A }_{oe,\rii} \\ {\bold A }_{oo,\rii}


\end{pmatrix}


,\qquad


{\bold A}_{ee,\rii} =


\begin{pmatrix}


{\bold B}_1 & {\bold O} & {\bold D} & {\bold O} & {\bold B}_1 &


{\bold O} & {\bold D} & {\bold O} \\


{\bold B}_2 & {\bold O} & {\bold O} & {\bold O} & {\bold B}_2 &


{\bold O} & {\bold O} & {\bold O} \\


{\bold B}_3 & {\bold O} & {\bold O} & {\bold O} & {\bold B}_3 &


{\bold O} & {\bold O} & {\bold O} \\


\vdots & \vdots & \vdots & \vdots & \vdots & \vdots & \vdots & \vdots \\


{\bold B}_N & {\bold O} & {\bold O} & {\bold O} & {\bold B}_N &


{\bold O} & {\bold O} & {\bold O}


\end{pmatrix}.


\end{equation*}


При этом блок ${\bold A}_{ee,\rii}$ состоит из диагональных матриц


$\bold D$ и матриц ${\bold B}_i$ размера $N\times N $, в каждой из


которых только один ненулевой коэффициент $({\bold B}_i)_{1,i}\ne 0$


при всех $i=1,\dotsc,N$. Структура матриц ${\bold A }_{eo,\rii}$, ${\bold A


}_{oe,\rii}$ и ${\bold A }_{oo,\rii}$ совпадает со структурой ${\bold A


}_{ee,\rii}$ с точностью до перестановки блочных столбцов


\cite{AGKS2002}. Следовательно, умножение матрицы ${\bold


A}_{\ri,\rii}$ на вектор требует


${\cal O}(p)$ арифметических действий, если умножение


производится только для ненулевых коэффициентов.





Общая вычислительная стоимость операции ${\bold v}\leftarrow\bi{\cal


K}_{\rii}^{(+)}{\bold r}$ складывается из стоимости решения системы с


матрицей $\bi{\cal S}$ (см. выше) и стоимости операций


продолжения и сужения. При этом последняя является наибольшей и


определяет главный член общей вычислительной стоимости. Она есть 


произведение числа итераций (\ref{inexact1}), которое при использовании 


любого из указанных солверов для внутренних задач Дирихле не превосходит


${\cal O}((1+\ln p)(1+\ln(1+\ln p)))$, и числа действий на одной итерации, 


требуемых для решения внутренних задач Дирихле.





\begin{thms}\label{T:gcost}


Пусть для решения локальных задач Дирихле применяется быстрый солвер 


c общей вычислительной стоимостью ${\mc O}(p^2(1+\ln p)^k)$, а 


предобусловливание дополнения Шура и продолжение с интерфейса


осуществляется согласно п.\,$1.4$. Пусть также $\Omega$ --- область,


составленная из прямоугольников, и сетка конечных элементов прямоугольная. 


Тогда общая вычислительная стоимость солвера МДО для системы 


$(\ref{slau})$ оценивается в


${\mc O}(p^2(1+\ln p)^{2{,}5+k}(1+\ln(1+\ln p)))$ арифметических действий.


\end{thms}





В численных экспериментах, представленных в разделе \ref{s:exper}, 


локальные задачи Дирихле решались посредством версий, описанных в п.\,1.3,


многоуровневого солвера и солвера типа МДО с крайними из указанных в


п.\,1.3 вычислительными стоимостями, соответствующими $k=0{,}5$ и $k=4{,}5$.


Заметим, что в общем случае (криволинейные конечные элементы, переменные


коэффициенты эллиптического уравнения) оценка теоремы \ref{T:gcost} 


сохраняется без учета стоимости операций умножения на каждой итерации 


на матрицу ${\bold K}$.








\section{Результаты численных экспериментов}\label{s:exper}


\label{IPCG}





2.1. {\it Общая характеристика}. Численные эксперименты производились 


с отдельными компонентами --- солверами МДО, и с МДО в целом.





Основной и, как правило, наиболее трудоемкой компонентой МДО является 


солвер для локальных задач Дирихле. Для солвера,


основанного на вторичном МДО \cite{K2001}, \cite{K2002}, который


используется при решении систем с предобусловливателем локальных


задач Дирихле, приводим часть результатов тестирования,


полученых А.\,Рытовым.


Результаты численных экспериментов с многосеточным солвером можно


найти в \cite{B2002} и \cite{KXA2002}. Практическая трудоемкость


обоих солверов в широком диапазоне значений $p$ подтверждает, что


они относятся к классу быстрых.





Довольно обширное экспериментирование производилось с операциями


продолжения с интерфейса посредством вторичных итерационных


процессов с предобусловливателями--солверами локальных задач


Дирихле. Здесь также подтвердился основной вывод из априорных


оценок, согласно которому для обеспечения сходимости внешнего


итерационного процесса и ее ускорения число итераций при


продолжении должно быть в ${\mc O}(\ln (1+p))$ раз больше, чем при


непосредственном решении локальных задач Дирихле.


Была также численно исследована зависимость обобщенного числа


обусловленности дополнения Шура при предобусловливании посредством


описанного в п.\,1.4 предобусловливателя.





Эффективность МДО в целом проверялась по отношению к


целому ряду его параметров. Решались задачи Дирихле в квадратной и


$L$-образной областях. При этом нас интересовали зависимости


числа глобальных итераций предобусловленного метода сопряженных


градиентов и времени счета от $p$ и количества вторичных итераций.


Наконец, сравнивались затраты процессорного времени по отношению к


погрешности решений, получаемых посредством описанного в работе


МДО на базе $p$-версии и многосеточного солвера коммерческого


пакета ANSYS для $h$-версии.


\medskip





2.2. {\it Решение локальных задач Дирихле}.


Приводим часть результатов, которые


касаются почти оптимального солвера, основанного на вторичной


декомпозиции \cite{K2001}, \cite{K2002}. Применялась наиболее


простая для программирования версия алгоритма, для которого


априорная оценка вычислительной стоимости имеет вид ${\mc O}(p^2(1+\ln


p)^{4{,}5})$, \cite{K2001},~\cite{K2002}. Время решения


(в~сек.) системы линейных алгебраических уравнений с


предобусловливателем в качестве матрицы и постоянной правой частью


в зависимости от $N$ ($p=2N+1$) приведено на рис.\,1.


Маркерами отмечены значения, соответствующие $\cal O (N^2 (\ln


N)^{3{,}5})$. Вычисления проводились на компьютере Pentium~IV 2.4\;GHz


с 1\;GB оперативной памяти. Как видно из графика, на практике солвер


оказывается более быстродействующим по сравнению с прогнозом на


основе априорной оценки.


\medskip





2.3. {\it Операции продолжения--сужения}.


Продолжения~(\ref{prolong}) и сужения~(\ref{restr}) выполнялись 


посредством вторичного итерационного процесса --- предобусловленного 


метода сопряженных градиентов с предобусловливателем~(\ref{Gbeu}) и 


выбором итерационного параметра $\beta_k$ в соответствии


с~(\ref{betak}). При этом варьировалось число вторичных итераций. 


Численные эксперименты позволили установить оптимальную (в смысле 


уменьшения вычислительных затрат глобального итерационного


солвера) зависимость количества итераций $n_{pr}$ вторичного итерационного 


солвера от $p$ (см. рис.\,2a)), которая полностью согласуется с 


теоретической оценкой $n_{\mathrm{pr}}(p)={\mc


O}((\ln (1+p))/\rho)$.





%\begin{figure}[htbp]


%\hspace{2cm}


%\special{em:graph ak-pic1.bmp}


%\vspace{5.5cm}


%\caption{Вычислительная стоимость солвера для локальных задач Дирихле.}


%\label{LDP}


%\end{figure}





\begin{center}


\unitlength=1mm


\begin{picture}(170,62) %% width, heigth, added 5 mm for signature


\put(45,62){\special{em:graph ak-pic1.bmp}} %% left X, upper Y


\put(18,0){\mbox Рис.\,1. Вычислительная стоимость солвера для 


локальных задач Дирихле}


\end{picture}


\end{center}








%\begin{figure}[htbp]


%\special{em:graph np_AND_N_63_2c.bmp}


%\caption{ a) оптимальное число итераций $n_{pr}$ вторичного итерационного солвера в зависимости от $p$;


%б) зависимость невязки приближенного решения от числа итераций глобального солвера.}


%\label{fignp}


%\end{figure}





\begin{center}


\unitlength=1mm


\begin{picture}(170,67) %% width, heigth, added 5 mm for signature


\put(2,67){\special{em:graph ak-pic2.bmp}} %% left X, upper Y


\put(2,5){\mbox Рис.\,2. a) оптимальное число итераций $n_{pr}$ 


вторичного итерационного солвера в зависимости}


\put(0,0){\mbox от $p$; б) зависимость невязки приближенного решения от числа итераций 


глобального солвера}


\end{picture}


\end{center}





Исследование влияния числа итераций $n_{pr}$ во вторичном солвере,


применяемом в продолжениях и сужениях, на эффективность глобального 


итерационного солвера позволяет сделать также следующие выводы. 


Уменьшение $n_{pr}$ приводит к потере сходимости внешнего итерационного 


процесса. При увеличении $n_{pr}$ во вторичном солвере число итераций 


$n_g$ внешнего процесса уменьшается, однако общие вычислительные затраты 


возрастают. На рис.\,2б) представлено уменьшение невязки в евклидовой 


норме $\gamma_k = \| {\bold K}{\bold u}^{(k)} -{\bold f}


\|_2/ \|{\bold f} \|_2$ при решении системы~(\ref{slau}) в зависимости 


от числа итераций глобального солвера $k$ для $p=127$ и различных 


значений $n_{pr}$. Аналогичные результаты получаются и для


других значений $p$.


\medskip





2.4. {\it Предобусловливание дополнения Шура}.


Решалась задача на собственные значения 


$\Bbb{S}_{\cal E} \bold x=\lambda \bold x$ и обобщенная


задача на собственные значения 


$\Bbb{S}_{\cal E}\bold x=\lambda\bi{\cal S}_{\cal E} \bold


x$. Для случая ${\bold A}={\bold A}_1$ зависимость числа обусловленности 


дополнения Шура от $p$ 


приведена на рис.\,3a), а рис.\,3б) содержит зависимость от $p$ обобщенного


числа обусловленности $\cond(\Bbb{S}_{\cal E}^{-1}{\bold S}_{\cal E})$. 


Графики приведены раздельно по причине существенной разницы в скоростях 


роста этих чисел обусловленности. Численные эксперименты полностью 


подтвердили априорные оценки теоремы~\ref{TeorSE}


асимптотического поведения обобщенного числа обусловленности.


\medskip





2.5. {\it МДО для задачи Дирихле в $L$-образной области}.


Для $L$-образной области, изображенной на рис.\,4, решалась задача Дирихле 


для уравнения Пуассона


\begin{equation}


\label{pois} u_{xx}'' + u_{yy}''=-f(x,y)


\end{equation}


с однородными краевыми условиями для $f(x,y)\equiv 1$. Система линейных 


алгебраических уравнений $p$-версии (\ref{slau}) решалась при помощи 


ПМСГ, описанного в п.\,1.2, с предобусловливателем МДО и итерационным 


параметром $\beta_k$, определенным по формуле (\ref{betak}).


Число $n_g$ итераций\linebreak ПМСГ, необходимых для уменьшения погрешности по 


невязке в евклидовой норме в $10^4$ раз, в зависимости от $p$ приведено на


рис.\,5а), а время решения системы $t$ (в~сек.) на компьютере 


Pentium~III-733\;Mhz c 768\,Mb оперативной памяти


--- на рис.\,5б). Поведение числа итераций и вычислительных затрат 


согласуется с априорными оценками.





%\begin{figure}[htpb]


%\special{em:graph shur_AND_precshur_2c.bmp}


%\caption{а) зависимость числа обусловленности дополнения Шура $\mbox{cond}(\mbb{S}_{\cal E})$


%от $p$; б) зависимость обобщенного числа обусловленности $\mbox{cond}(\bi{\cal S}_{\cal


% E}^{-1}\mbb{S}_{\cal E})$ от $p$.}


%\label{figshur}


%\end{figure}





\begin{center}


\unitlength=1mm


\begin{picture}(170,67) %% width, heigth, added 5 mm for signature


\put(5,67){\special {em:graph ak-pic3.bmp}} %%left X, upper Y


\put(15,5){\mbox Рис.\,3. а) зависимость числа обусловленности дополнения 


Шура $\cond(\Bbb{S}_{\cal E})$ от $p$;}


\put(23,0){\mbox б) зависимость обобщенного числа 


обусловленности $\cond(\bi{\cal S}_{\cal


E}^{-1}\Bbb{S}_{\cal E})$ от $p$}


\end{picture}


\end{center}








%\begin{figure}[h!]


%\hspace{6cm}


%\special{em:graph Lshape_2c.bmp}


%\vspace{4cm}


%\caption{$L$--образная область.}


%\label{figLSD}


%\end{figure}





\begin{center}


\unitlength=1mm


\begin{picture}(170,45) %% width, heigth, added 5 mm for signature


\put(65,45){\special {em:graph ak-pic4.bmp}} %% left X, upper Y,


\put(60,0){\mbox Рис.\,4. $L$-образная область}


\end{picture}


\end{center}








%\begin{figure}[h!]


%\hspace{0.5cm}


%\special{em:graph LIterP_AND_LIterT_2c.bmp}


%\vspace{5.6cm}


%\caption{а) зависимость числа итераций ПМСГ от $p$;


%б) зависимость времени решения системы линейных уравнений $p$--версии от $p$.}


%\label{figLIterP}


%\end{figure}





\begin{center}


\unitlength=1mm


\begin{picture}(170,67) %% width, heigth, added 5 mm for signature


\put(5,67){\special {em:graph ak-pic5.bmp}} %% left X, upper Y


\put(40,5){\mbox Рис.\,5. а) зависимость числа итераций ПМСГ от $p$;}


\put(16,0){\mbox б) зависимость времени решения системы линейных уравнений


$p$-версии от $p$}


\end{picture}


\end{center}








2.6. {\it $p$- и $h$-версии в задаче с гладким решением}.


В случае гладких решений скорость сходимости $p$-версии, как правило, 


выше скорости сходимости $h$-версии (напр., \cite{BS1987},


\cite{Schwab1998}). Представляет интерес сравнение времени


решения систем линейных алгебраических уравнений для $h$- и 


$p$-дис\-кретизаций, необходимых для получения приближенного решения 


с одинаковой относительной погрешностью в энергетической норме. В


нашем численном эксперименте решалась однородная граничная задача 


Дирихле для уравнения Пуассона~(\ref{pois}) в квадрате $(-4,4)\times(4,4)$


c правой частью $f(x,y)=2\pi^2\sin\pi x\sin\pi y$, соответствующей 


точному решению $u_{ex}(x,y)=\sin\pi x\sin\pi y$.





Временные затраты $t$ (в сек.) на решение систем линейных уравнений 


в зависимости от числа степеней свободы $N_{DOF}$ приведены на 


рис.\,6a) (для $h$-версии) и рис.~6б) (для


$p$-версии). Отметим, что для $h$-версии $N_{DOF}={\cal O}(h^{-2})$,


а в случае $p$-версии $N_{DOF}={\cal O}(p^{2})$. Решение системы 


линейных уравнений $h$-версии осуществлялось в ANSYS при помощи 


многосеточного солвера (Algebraic multigrid solver). В нашем пакете 


программ реализован итерационный солвер с МДО-предобусловливателем, 


описанный в пп.\,1.2--1.4. Рисунок~7a) содержит сравнение вычислительных 


затрат $t$ (в~сек.) прямого и итерационного солверов для решения систем 


линейных уравнений, соответственно $h$- и $p$-дискретизаций, в зависимости 


от относительной ошибки приближенного решения в энергетической норме.


%\begin{figure}[h!]


% \hspace{0.5cm}


% \special{em:graph hTime_AND_pTime_2c.bmp}


% \vspace{5.6cm}


% \caption{а) время решения системы линейных уравнений $h$--версии в зависимости от числа степеней свободы;


% б) время решения системы линейных уравнений $p$--версии в зависимости от числа степеней


% свободы.}


% \label{fighTime}


% \end{figure}


% \begin{figure}[h!]


% \hspace{0.5cm}


% \special{em:graph hpEnTime_AND_energy_2c.bmp}


% \vspace{5.6cm}


% \caption{а) сопоставление вычислительных затрат на решение систем линейных уравнений $p$ и $h$--версий;


%б) зависимость относительной погрешности конечноэлементного решения в энергетической норме от числа


%степеней свободы $N_{DOF}$.}


% \label{fighpEnTime}


%\end{figure}





\begin{center}


\unitlength=1mm


\begin{picture}(170,72) %% width, heigth, added 5 mm for signature


\put(2,72){\special {em:graph ak-pic6.bmp}} %% left X, upper Y


\put(5,10){\mbox Рис.\,6. а) время решения системы линейных уравнений


$h$-версии в зависимости от числа}


\put(4,5){\mbox степеней свободы; б) время решения 


системы линейных уравнений $p$-версии в зависимости}


\put(60,0){\mbox от числа степеней свободы}


\end{picture}


\end{center}





\begin{center}


\unitlength=1mm


\begin{picture}(170,72) %% width, heigth, added 5 mm for signature


\put(2,72){\special {em:graph ak-pic7.bmp}} %% left X, upper Y


\put(4,10){\mbox Рис.\,7. а) сопоставление вычислительных затрат на 


решение систем линейных уравнений $p$-}


\put(10,5){\mbox и $h$-версий; б) зависимость 


относительной погрешности конечноэлементного решения}


\put(30,0){\mbox в энергетической 


норме от числа степеней свободы $N_{DOF}$}


\end{picture}


\end{center}








Исследовалось также уменьшение относительной погрешности в энергетической 


норме $\delta_E$ при увеличении числа степеней свобод $N_{DOF}$. Решение 


при помощи $h$-версии было получено в пакете ANSYS, а для $p$-версии 


использовался наш пакет. Полученные результаты (см. рис.\,7б)) полностью 


согласуются с априорными оценками скорости сходимости $h$- и


$p$-версий метода конечных элементов.








\begin{thebibliography}{99}





\bibitem{BS1987} %1


Babu\v{s}ka I., Suri M. {\it The $h$-$p$ version of the finite element 


method with quasiuniform meshes} // Math. Mod. Numer. Anal. --


1987. -- V.\,21. -- P.\,199--238.





\bibitem{Schwab1998} %2


Schwab Ch. {\it $P$- and $hp$-finite element methods. Theory and 


applications in solid and fluid mechanics}. -- Clarendon Press. Oxford,


1998. -- 374~p.





\bibitem{Nepomnyaschikh:91} %3


Nepomnyaschikh S.V. {\it Method of splitting into subspaces for solving 


elliptic boundary value problems in complex-form domains} //


Russian (cont. Sov.) J. Numer. Anal. and


Math. Model. -- 1991. -- V.\,6. -- \No\,2. -- P.\,151--168.





\bibitem{SmithBjorstadGropp:96} %4


Smith B., Bj{\o}rstad P., Gropp W. {\it Domain decomposition\/${:}$


parallel multilevel methods for elliptic partial differential 


equations}. -- Cambride University Press, 1996. -- 224 p. %% (+xii)





\bibitem{BabuskaCraigMandelPitkaranta:91} %5


Babu\v{s}ka I., Craig A., Mandel J., Pitk\"{a}ranta J. {\it Efficient 


preconditioning for the $p$-version finite element method in two 


dimensions} // SIAM J. Numer. Anal. -- 1991. -- V.\,28. -- \No\,3.


-- P.\,624--661.





\bibitem{Pavarino:94} %6


Pavarino L.F. {\it Additive Schwartz methods for the $p$-version finite


element method} // Numer. Math. -- 1994. -- V.\,66. -- \No\,4. 


-- P.\,493--515.





\bibitem{ainsworth} %7


Ainsworth M.


{\it A preconditioner based on domain decomposition for $h$-$p$


finite element approximation on quasi-uniform meshes} //


SIAM J. Numer. Anal. -- 1996. -- V.\,33. -- P.\,1358--1376.





\bibitem{AinsworthSenior} %8


Ainsworth M., Senior B.


{\it Aspects of an adaptive $hp$-finite element method\/${:}$ Adaptive


strategy, conforming approximation and efficient solvers} //


Comp. Methods Appl. Mech. Engrg. -- 1997. -- V.\,150. -- P.\,65--87.





\bibitem{guo} %9


Guo B., Cao W. {\it Inexact solvers on element surfaces for the $p$


and $h$-$p$ finite element method} // Comput.


Methods Appl. Mech. -- 1997. -- V.\,150. -- P.\,173--191.





\bibitem{KI1995} %10


Ivanov S.A., Korneev V.G. {\it On the preconditioning in the domain 


decomposition technique for the $p$-version finite element method}.


I // Technishe Universit\"{a}t


Chemnitz-Zwickau. Preprint SPC 95-35. -- 1995. -- \No\,12. -- P.\,1--15.





\bibitem{KI1995a} %11


Ivanov S.A., Korneev V.G. {\it On the preconditioning in the domain 


decomposition technique for the $p$-version finite element method}.


II // Technishe Universit\"{a}t


Chemnitz-Zwickau. Preprint SPC 95-36. -- 1995. -- \No\,12. -- P.\,1--14.





\bibitem{KI1996} %12


Ivanov S.A., Korneev V.G. {\it Preconditioning in the domain decomposition


methods for the $p$-version with the hierarchical bases} //


Матем. моделир. -- 1996. -- Т.\,8. -- \No\,9. -- С.\,63--73.





\bibitem{KJ1997} %13


Korneev V.G., Jensen S. {\it On domain decomposition preconditioning in 


the hierarchical $p$-version of the finite element method} // Comput. 


Methods. Appl. Mech. Eng. -- 1997. -- V.\,150. -- P.\,215--238.





\bibitem{KJ1999} %14


Корнеев В.Г., Енсен С. {\it Эффективное предобусловливание методом 


декомпозиции области для $p$-версии с


иерархическим базисом}. I // Изв. вузов. Математика. --


1999. -- \No\,5. -- C.\,37--56.





\bibitem{KJ1999a} %15


Корнеев В.Г., Енсен С. {\it Эффективное предобусловливание методом 


декомпозиции области для $p$-версии с


иерархическим базисом}. II // Изв. вузов. Математика. --


1999. -- \No\,11. -- C.\,24--40.





\bibitem{K2001} %16


Корнеев В.Г. {\it Почти оптимальный метод решения задач Дирихле на 


подобластях декомпозиции иерархической $hp$-версии} // Дифференц. 


уравнения. -- 2001. -- Т.\,37. -- \No\,7. -- С.\,1--10.





\bibitem{K2002} %17


Korneev V.G. {\it Local Dirichlet problems on subdomains of decomposition 


in $hp$ discretizations, and optimal algorithms for their solution} //


Матем. моделир. -- 2002. -- Т.14. -- \No\,5. -- С.\,51--74.





\bibitem{Korneev:02b} %18


Korneev V.G. {\it Fast domain decomposition solvers for $hp$ 


discretizations of $2$-nd order elliptic equations} // In Proceedings 


of the International Conference


on Computational Mathematics. 24--28 of June, 2002, Novosibirsk. Institute 


of Computational Mechanics and Mathematical Geophysics of


Sibirian Branch of Russian Acad. of Sci.: Novosibirsk, 2002. -- P.\,536--547.





\bibitem{KXA2002} %19


Korneev V.G., Xanthis L., Anoufriev I.E.


{\it Solving finite element $hp$ discretizations of elliptic problems by fast 


DD algorithms} // Тр. СПбГПУ. -- 2002. -- \No\,485. -- С.\,126--153.





\bibitem{K1977} %20


Корнеев В.Г. {\it Схемы метода конечных элементов высоких порядков 


точности}. -- Л.: Изд-во ЛГУ, 1977. -- 205~с.





\bibitem{GGK} %21


Giladi E., Golub G., Keller J. {\it Inner and outer iterations for the 


Chebyshev algorithm} // SIAM J. on Numer. Anal. -- 1998. --


V.\,35. -- \No\,1. -- P.\,300--319.





\bibitem{GE} %22


Golub G.H., Ye Q. {\it Inexact preconditioned conjugate gradient method 


with inner--outer iteration} // SIAM J. Sc. Comp. -- 1999. --


V.\,21. -- \No\,4. -- P.\,1305--1320.





\bibitem{GO} %23


Golub G., Overton M. {\it The convergence of inexact Chebyshev and 


Richardson iterative methods for solving linear systems} // Numer. Math. 


-- 1988. -- V.\,53. -- P.\,571--593.





\bibitem{B2002} %24


Beuchler S. {\it Multigrid solver for the inner problem in domain 


decomposition methods for $p$-FEM} //


SIAM J. Numer. Anal. -- 2002. -- V.\,40. -- \No\,4. -- P.\,928--944.





\bibitem{BSS2002} %25


Beuchler S., Schneider R., Schwab C.


{\it Multiresolution weighted norm equivalence and applications} //


Preprint SFB393/02-09, Technische Universit\"at Chemnitz, 2002.





\bibitem{IvanovKorneev:95} %26


Иванов С.А., Корнеев В.Г. {\it Выбор координатных функций высоких 


порядков и предобусловливание в рамках метода декомпозиции области} //


Изв. вузов. Mатематика. -- 1995. -- \No\,4. -- С.\,62--81.





\bibitem{KXA2002P} %27


Korneev V., Xantis L., Anoufriev I. {\it Hierarchical and Lagrange $hp$ 


discretizations and fast domain decomposition solvers for them} //


Johannes Kepler Universit\"{a}t Linz. SFB-Report


\No\,02-18, 2002. -- 30~p.





\bibitem{AGKS2002} %28


Ануфриев И.Е., Головин О.А., Константинов В., Смирнова Е.Н.


{\it Численный эксперимент с эффективным предобусловливанием интерфейсной 


задачи в $p$-версии} //


Тр. СПбГПУ. -- 2002. -- \No\,485. -- С.\,3--16.








\end{thebibliography}





\vskip 2cm





\post{\it Санкт-Петербургский}{\it Поступила}


\post{\it государственный политехнический}{20.10.2003}


\postup{\it университет}{} 


\label{end}


\end{document}








%%%%%% \bibitem{Rytov} А.~Рытов. СТАТЬЯ ПРО ЧИСЛЕННЫЕ ЭКСПЕРИМЕНТЫ. 2003.





